\documentclass[10pt,a4paper]{amsart}
\usepackage{amsmath,amssymb,amsthm,graphicx}
\usepackage{verbatim}
\newenvironment{reference}{\comment}{\endcomment}
%\newenvironment{reference}{}{}
\newenvironment{slogan}{\comment}{\endcomment}
\newenvironment{history}{\comment}{\endcomment}

% For commutative diagrams we use Xy-pic
\usepackage[all]{xy}

% We use 2cell for 2-commutative diagrams.
\xyoption{2cell}
% Theorem environments.
%
\theoremstyle{plain}
\newtheorem{theorem}[subsection]{Theorem}
\newtheorem{proposition}[subsection]{Proposition}
\newtheorem{lemma}[subsection]{Lemma}

\theoremstyle{definition}
\newtheorem{definition}[subsection]{Definition}
\newtheorem{example}[subsection]{Example}
\newtheorem{exercise}[subsection]{Exercise}
\newtheorem{situation}[subsection]{Situation}

\theoremstyle{remark}
\newtheorem{remark}[subsection]{Remark}
\newtheorem{remarks}[subsection]{Remarks}

\numberwithin{equation}{subsection}

% Macros
%
\def\lim{\mathop{\mathrm{lim}}\nolimits}
\def\colim{\mathop{\mathrm{colim}}\nolimits}
\def\Spec{\mathop{\mathrm{Spec}}}
\def\Hom{\mathop{\mathrm{Hom}}\nolimits}
\def\Ext{\mathop{\mathrm{Ext}}\nolimits}
\def\SheafHom{\mathop{\mathcal{H}\!\mathit{om}}\nolimits}
\def\SheafExt{\mathop{\mathcal{E}\!\mathit{xt}}\nolimits}
\def\Sch{\mathit{Sch}}
\def\Mor{\mathop{\mathrm{Mor}}\nolimits}
\def\Ob{\mathop{\mathrm{Ob}}\nolimits}
\def\Sh{\mathop{\mathit{Sh}}\nolimits}
\def\NL{\mathop{N\!L}\nolimits}
\def\CH{\mathop{\mathrm{CH}}\nolimits}
\def\proetale{{pro\text{-}\acute{e}tale}}
\def\etale{{\acute{e}tale}}
\def\QCoh{\mathit{QCoh}}
\def\Ker{\mathop{\mathrm{Ker}}}
\def\Im{\mathop{\mathrm{Im}}}
\def\Coker{\mathop{\mathrm{Coker}}}
\def\Coim{\mathop{\mathrm{Coim}}}

% Boxtimes
%
\DeclareMathSymbol{\boxtimes}{\mathbin}{AMSa}{"02}

%
% Macros for moduli stacks/spaces
%
\def\QCohstack{\mathcal{QC}\!\mathit{oh}}
\def\Cohstack{\mathcal{C}\!\mathit{oh}}
\def\Spacesstack{\mathcal{S}\!\mathit{paces}}
\def\Quotfunctor{\mathrm{Quot}}
\def\Hilbfunctor{\mathrm{Hilb}}
\def\Curvesstack{\mathcal{C}\!\mathit{urves}}
\def\Polarizedstack{\mathcal{P}\!\mathit{olarized}}
\def\Complexesstack{\mathcal{C}\!\mathit{omplexes}}
% \Pic is the operator that assigns to X its picard group, usage \Pic(X)
% \Picardstack_{X/B} denotes the Picard stack of X over B
% \Picardfunctor_{X/B} denotes the Picard functor of X over B
\def\Pic{\mathop{\mathrm{Pic}}\nolimits}
\def\Picardstack{\mathcal{P}\!\mathit{ic}}
\def\Picardfunctor{\mathrm{Pic}}
\def\Deformationcategory{\mathcal{D}\!\mathit{ef}}
\usepackage{cleveref}
\crefname{originalreference}{original reference}{original references}
\Crefname{originalreference}{Original reference}{Original references}
\makeatletter
\newlabel{definition-big-proetale-site}{{098G}{source}{}{sourcepointer-1}{}}
\newlabel{definition-big-proetale-site@cref}{{[originalreference][0][]098G}{[1][0][]source}{}{}{}}
\makeatother
\crefname{originalreference}{original reference}{original references}
\Crefname{originalreference}{Original reference}{Original references}
\makeatletter
\newlabel{lemma-etale-proetale}{{098B}{source}{}{sourcepointer-2}{}}
\newlabel{lemma-etale-proetale@cref}{{[originalreference][0][]098B}{[1][0][]source}{}{}{}}
\makeatother
\crefname{originalreference}{original reference}{original references}
\Crefname{originalreference}{Original reference}{Original references}
\makeatletter
\newlabel{lemma-proetale-subcanonical}{{098Z}{source}{}{sourcepointer-3}{}}
\newlabel{lemma-proetale-subcanonical@cref}{{[originalreference][0][]098Z}{[1][0][]source}{}{}{}}
\makeatother
\crefname{originalreference}{original reference}{original references}
\Crefname{originalreference}{Original reference}{Original references}
\makeatletter
\newlabel{lemma-w-contractible-proetale-cover}{{098F}{source}{}{sourcepointer-4}{}}
\newlabel{lemma-w-contractible-proetale-cover@cref}{{[originalreference][0][]098F}{[1][0][]source}{}{}{}}
\makeatother
\crefname{originalreference}{original reference}{original references}
\Crefname{originalreference}{Original reference}{Original references}
\makeatletter
\newlabel{proposition-find-w-contractible}{{0983}{source}{}{sourcepointer-5}{}}
\newlabel{proposition-find-w-contractible@cref}{{[originalreference][0][]0983}{[1][0][]source}{}{}{}}
\makeatother
\crefname{originalreference}{original reference}{original references}
\Crefname{originalreference}{Original reference}{Original references}
\makeatletter
\newlabel{remark-size-w-contractible}{{0984}{source}{}{sourcepointer-6}{}}
\newlabel{remark-size-w-contractible@cref}{{[originalreference][0][]0984}{[1][0][]source}{}{}{}}
\makeatother
\crefname{originallemma}{original lemma}{original lemmas}
\Crefname{originallemma}{Original lemma}{Original lemmas}
\makeatletter
\newlabel{sets-lemma-bound-size}{{000Q}{source}{}{sourcepointer-7}{}}
\newlabel{sets-lemma-bound-size@cref}{{[originallemma][0][]000Q}{[1][0][]source}{}{}{}}
\makeatother
\crefname{originaldefinition}{original definition}{original definitions}
\Crefname{originaldefinition}{Original definition}{Original definitions}
\makeatletter
\newlabel{sites-definition-sheaves-injective-surjective}{{00WM}{source}{}{sourcepointer-8}{}}
\newlabel{sites-definition-sheaves-injective-surjective@cref}{{[originaldefinition][0][]00WM}{[1][0][]source}{}{}{}}
\makeatother
\crefname{originaldefinition}{original definition}{original definitions}
\Crefname{originaldefinition}{Original definition}{Original definitions}
\makeatletter
\newlabel{sites-definition-w-contractible}{{090L}{source}{}{sourcepointer-9}{}}
\newlabel{sites-definition-w-contractible@cref}{{[originaldefinition][0][]090L}{[1][0][]source}{}{}{}}
\makeatother
\crefname{originallemma}{original lemma}{original lemmas}
\Crefname{originallemma}{Original lemma}{Original lemmas}
\makeatletter
\newlabel{sites-lemma-covering-surjective-after-sheafification}{{00WT}{source}{}{sourcepointer-10}{}}
\newlabel{sites-lemma-covering-surjective-after-sheafification@cref}{{[originallemma][0][]00WT}{[1][0][]source}{}{}{}}
\makeatother
\title{Characterization of w-contractible affine schemes as weakly contractible objects of the big pro-etale site}
\author{The Stacks Project Authors}
\date{stacks-project:proetale.tex:local-snapshot-20261012; GNU FDL 1.2 or later}
\begin{document}
\maketitle
\noindent\textit{Editorial note.} All delivered mathematical text is the author's own, copied from the pinned Stacks Project chapter file proetale.tex; no statement or proof step was rewritten, reordered or removed.\par\medskip
\noindent\textit{Editorial note.} The counted proof is the author's complete argument: the deduction of weak contractibility from the refinement property through finite Zariski covers and the induced product decompositions, and the deduction of w-contractibility by turning a faithfully flat weakly etale extension into a covering and producing the retraction from weak contractibility.\par\medskip
\noindent\textit{Editorial note.} The label lemma-w-contractible-is-weakly-contractible names the counted statement, and the proof environment immediately following it in the source is delivered verbatim as its proof.\par\medskip
\noindent\textit{Editorial note.} The counted statement is carried by the author's own lemma environment in the chapter source.\par\medskip
\noindent\textit{Editorial note.} Pre-open gate: python operations/agent-authoring/tools/precard\_gate.py --source proetale.tex --label lemma-w-contractible-is-weakly-contractible --start 3026 --end 3084 --local-id 3582 returned RESULT=OK (L3582 free, occupancy FREE, public interval clear). Candidate taken from sources/stacks-proetale/candidates-halfA.json.\par\medskip
\noindent\textit{Editorial note.} The statement and proof are copied from the pinned chapter file; the cited inputs are kept as external pointers with their canonical Stacks tags.\par\medskip
\noindent\textit{Allowed standard prerequisites (editorial).}\begin{itemize}
\item Equivalence of the ring theoretic condition with refinement of pro-etale coverings by finite Zariski coverings. Same chapter, the lemma identifying w-contractibility of an affine scheme with the property that every pro-etale covering is refined by a finite Zariski covering. Supplies the equivalence of parts (1) and (3).
\item Construction of a faithfully flat weakly etale extension with a retraction. Same chapter, the proposition producing the extension A -> D of an affine w-local-like situation together with the retraction results used at the end. Supplies the object V = Spec(D) covering T used to produce the retraction sigma and hence the retractions of faithfully flat weakly etale extensions.
\item Subcanonicity of the pro-etale topology and surjectivity of the representable maps. Same chapter, the lemma that the pro-etale topology is subcanonical, and Sites, the lemma turning a covering into a surjective map of sheaves after sheafification. Supply the surjectivity of h\_V -> h\_T, which combines with weak contractibility to produce a section of T over V.
\item The size of affine schemes and the bound on the big pro-etale site. Sets, the lemma computing the size of affine schemes, together with the definition of the big pro-etale site and the size estimate for the weakly contractible ring. Supplies the membership of V as an object of the big pro-etale site.
\end{itemize}\medskip
\ifdefined\hypertarget\hypertarget{sourcepointer-1}{}\fi
\noindent\textit{Source reference.} Original reference 098G: sources/stacks-proetale/pinned-original-source.tex, source lines 2465--2465. Prerequisite pointer invoked at proof line(s) 3029; it names a standard fact of the same chapter that the counted argument uses, and every step of the counted proof built on it is retained verbatim in the delivered body.\par\medskip
\ifdefined\hypertarget\hypertarget{sourcepointer-2}{}\fi
\noindent\textit{Source reference.} Original reference 098B: sources/stacks-proetale/pinned-original-source.tex, source lines 2393--2393. Prerequisite pointer invoked at proof line(s) 3057; it names a standard fact of the same chapter that the counted argument uses, and every step of the counted proof built on it is retained verbatim in the delivered body.\par\medskip
\ifdefined\hypertarget\hypertarget{sourcepointer-3}{}\fi
\noindent\textit{Source reference.} Original reference 098Z: sources/stacks-proetale/pinned-original-source.tex, source lines 2966--2966. Prerequisite pointer invoked at proof line(s) 3073; it names a standard fact of the same chapter that the counted argument uses, and every step of the counted proof built on it is retained verbatim in the delivered body.\par\medskip
\ifdefined\hypertarget\hypertarget{sourcepointer-4}{}\fi
\noindent\textit{Source reference.} Original reference 098F: sources/stacks-proetale/pinned-original-source.tex, source lines 2998--2998. Prerequisite pointer invoked at proof line(s) 3044; it names a standard fact of the same chapter that the counted argument uses, and every step of the counted proof built on it is retained verbatim in the delivered body.\par\medskip
\ifdefined\hypertarget\hypertarget{sourcepointer-5}{}\fi
\noindent\textit{Source reference.} Original reference 0983: sources/stacks-proetale/pinned-original-source.tex, source lines 2135--2135. Prerequisite pointer invoked at proof line(s) 3064; it names a standard fact of the same chapter that the counted argument uses, and every step of the counted proof built on it is retained verbatim in the delivered body.\par\medskip
\ifdefined\hypertarget\hypertarget{sourcepointer-6}{}\fi
\noindent\textit{Source reference.} Original reference 0984: sources/stacks-proetale/pinned-original-source.tex, source lines 2171--2171. Prerequisite pointer invoked at proof line(s) 3067; it names a standard fact of the same chapter that the counted argument uses, and every step of the counted proof built on it is retained verbatim in the delivered body.\par\medskip
\ifdefined\hypertarget\hypertarget{sourcepointer-7}{}\fi
\noindent\textit{Source reference.} Original lemma 000Q: sources/stacks-chow/tags.txt, source lines 35--35. Prerequisite pointer invoked at proof line(s) 3071; it names a standard fact of a sibling chapter that the counted argument uses, and every step of the counted proof built on it is retained verbatim in the delivered body.\par\medskip
\ifdefined\hypertarget\hypertarget{sourcepointer-8}{}\fi
\noindent\textit{Source reference.} Original definition 00WM: sources/stacks-chow/tags.txt, source lines 1188--1188. Prerequisite pointer invoked at proof line(s) 3052; it names a standard fact of a sibling chapter that the counted argument uses, and every step of the counted proof built on it is retained verbatim in the delivered body.\par\medskip
\ifdefined\hypertarget\hypertarget{sourcepointer-9}{}\fi
\noindent\textit{Source reference.} Original definition 090L: sources/stacks-chow/tags.txt, source lines 11414--11414. Prerequisite pointer invoked at proof line(s) 3035; it names a standard fact of a sibling chapter that the counted argument uses, and every step of the counted proof built on it is retained verbatim in the delivered body.\par\medskip
\ifdefined\hypertarget\hypertarget{sourcepointer-10}{}\fi
\noindent\textit{Source reference.} Original lemma 00WT: sources/stacks-chow/tags.txt, source lines 1195--1195. Prerequisite pointer invoked at proof line(s) 3075; it names a standard fact of a sibling chapter that the counted argument uses, and every step of the counted proof built on it is retained verbatim in the delivered body.\par\medskip
\begin{lemma}
\label{lemma-w-contractible-is-weakly-contractible}
Let $\Sch_\proetale$ be a big pro-\'etale site as in
Definition \ref{definition-big-proetale-site}.
Let $T = \Spec(A)$ be an affine object of $\Sch_\proetale$.
The following are equivalent
\begin{enumerate}
\item $A$ is w-contractible,
\item $T$ is a weakly contractible
(Sites, Definition \ref{sites-definition-w-contractible})
object of $\Sch_\proetale$, and
\item every pro-\'etale covering of $T$ can be refined by
a Zariski covering of the form $T = \coprod_{i = 1, \ldots, n} U_i$.
\end{enumerate}
\end{lemma}

\begin{proof}
We have seen the equivalence of (1) and (3) in
Lemma \ref{lemma-w-contractible-proetale-cover}.

\medskip\noindent
Assume (3) and let $\mathcal{F} \to \mathcal{G}$ be a surjection of sheaves on
$\Sch_\proetale$. Let $s \in \mathcal{G}(T)$. To prove (2) we will show that
$s$ is in the image of $\mathcal{F}(T) \to \mathcal{G}(T)$. We can find a
covering $\{T_i \to T\}$ of $\Sch_\proetale$ such that $s$ lifts
to a section of $\mathcal{F}$ over $T_i$
(Sites, Definition \ref{sites-definition-sheaves-injective-surjective}).
By (3) we may assume we have a finite covering
$T = \coprod_{j = 1, \ldots, m} U_j$ by open and closed subsets
and we have $t_j \in \mathcal{F}(U_j)$ mapping to $s|_{U_j}$.
Since Zariski coverings are coverings in $\Sch_\proetale$
(Lemma \ref{lemma-etale-proetale}) we conclude that
$\mathcal{F}(T) = \prod \mathcal{F}(U_j)$.
Thus $t = (t_1, \ldots, t_m) \in \mathcal{F}(T)$
is a section mapping to $s$.

\medskip\noindent
Assume (2). Let $A \to D$ be as in
Proposition \ref{proposition-find-w-contractible}.
Then $\{V \to T\}$ is a covering of $\Sch_\proetale$.
(Note that $V = \Spec(D)$ is an object of $\Sch_\proetale$
by Remark \ref{remark-size-w-contractible}
combined with our choice of the function
$Bound$ in Definition \ref{definition-big-proetale-site}
and the computation of the size of affine schemes in
Sets, Lemma \ref{sets-lemma-bound-size}.)
Since the topology on $\Sch_\proetale$ is subcanonical
(Lemma \ref{lemma-proetale-subcanonical})
we see that $h_V \to h_T$ is a surjective map of sheaves
(Sites, Lemma \ref{sites-lemma-covering-surjective-after-sheafification}).
Since $T$ is assumed weakly contractible, we see that there is an element
$f \in h_V(T) = \Mor(T, V)$ whose image in $h_T(T)$ is $\text{id}_T$.
Thus $A \to D$ has a retraction $\sigma : D \to A$.
Now if $A \to B$ is faithfully flat and weakly \'etale, then
$D \to D \otimes_A B$ has the same properties, hence
there is a retraction $D \otimes_A B \to D$ and combined
with $\sigma$ we get a retraction $B \to D \otimes_A B \to D \to A$
of $A \to B$. Thus $A$ is w-contractible and (1) holds.
\end{proof}

\end{document}
